\documentclass[aspectratio=169]{pres}

\title[WeatherScope]{WeatherScope}
\subtitle{Observatoire météo historique sur NOAA GHCN-Daily\\Allemagne et Espagne, local puis distribué}
\author[Claveille, Jules, Diallo]{CLAVEILLE Hugo, JULES Stevenson, DIALLO Mamadou Ougailou\\[0.3em]\small M2 SLCD -- Big Data and NoSQL}
\date{10 novembre 2026}

% Résultats à compléter (distribué complet, E1-E4)
\newcommand{\afaire}[1]{\textcolor{red}{\textbf{[À compléter : #1]}}}
\lstdefinestyle{sql}{language=SQL, basicstyle=\ttfamily\scriptsize, keywordstyle=\color{accent}\bfseries,
  commentstyle=\color{gray}\itshape, morekeywords={FILTER,OVER,PARTITION,CONFLICT,DISTINCT,ILIKE}}

\begin{document}

\begin{frame}
  \titlepage
\end{frame}

\begin{frame}{Plan}
  \setcounter{tocdepth}{1}
  \tableofcontents
\end{frame}

%====================================================
\section{Périmètre et données}

\begin{frame}{Le périmètre}
  \begin{columns}[T]
    \begin{column}{0.5\textwidth}
      \begin{block}{Ce qu'on a enregistré}
        \begin{itemize}
          \item \textbf{Zones :} Allemagne (\code{GM}) et Espagne (\code{SP}), préfixe pays de l'identifiant NOAA
          \item \textbf{Étude :} 2015--2024
          \item \textbf{Référence :} 1991--2020 (normale sur 30 ans)
          \item \textbf{Éléments :} TMAX, TMIN, PRCP
        \end{itemize}
      \end{block}
    \end{column}
    \begin{column}{0.46\textwidth}
      \begin{block}{Volume chargé (1991--2024)}
        \begin{itemize}
          \item 774 stations (589 GM, 185 SP)
          \item \textbf{21,46 M} observations réelles
          \item 2,9 Go en base
        \end{itemize}
      \end{block}
      \small
      \textbf{Pourquoi :} deux climats contrastés ; l'étude seule (7,14 M) n'atteint pas les 10 M d'E1, la référence est nécessaire à F3 et F4.
    \end{column}
  \end{columns}
\end{frame}

\begin{frame}{Provenance et qualité}
  \begin{columns}[T]
    \begin{column}{0.48\textwidth}
      \begin{block}{Provenance}
        \begin{itemize}
          \item Manifeste versionné : URL, date, taille, \textbf{SHA-256} de chaque fichier
          \item Table \code{batches} : un lot par import
          \item Chaque observation garde son lot
        \end{itemize}
      \end{block}
      \begin{block}{Règle de validité}
        \code{qflag} vide et \code{mflag} $\neq$ P\\[0.2em]
        \small Trace (T) = 0 mm, comptée à part.\\
        Invalide : stockée, comptée, jamais utilisée.
      \end{block}
    \end{column}
    \begin{column}{0.48\textwidth}
      \setbeamercolor{block title}{fg=white,bg=red!80!black}
      \begin{block}{Anomalies trouvées}
        \small
        \begin{itemize}
          \item 1\,403 observations avec \code{qflag}, 267 traces, aucune P
          \item Heure d'observation vide pour 100\,\% des lignes
          \item 98\,\% des valeurs viennent d'ECA\&D (\code{sflag} E)
          \item Espagne 2024 : environ 24\,\% d'observations en moins
          \item 9 stations espagnoles avec une longitude de signe faux : documentées, pas corrigées
        \end{itemize}
      \end{block}
    \end{column}
  \end{columns}
\end{frame}

%====================================================
\section{Architecture et modèle}

\begin{frame}{Choix technologique : PostgreSQL, puis Citus}
  \begin{columns}[T]
    \begin{column}{0.52\textwidth}
      \begin{block}{Pourquoi}
        \begin{itemize}
          \item Données tabulaires de forme fixe
          \item F2--F4 en quelques lignes avec les fonctions de fenêtre SQL
          \item \textbf{Le même SQL} en local et sur Citus : seul \code{DATABASE\_URL} change
          \item \code{EXPLAIN} montre les workers contactés
          \item Échec explicite si un worker manque
        \end{itemize}
      \end{block}
    \end{column}
    \begin{column}{0.44\textwidth}
      \setbeamercolor{block title}{fg=white,bg=gray}
      \begin{block}{Écartés}
        \small
        \begin{itemize}
          \item \textbf{MongoDB :} F2--F4 en longs pipelines d'agrégation, trop risqué dans le délai
          \item \textbf{Cassandra, ClickHouse :} jamais pratiqués
        \end{itemize}
      \end{block}
    \end{column}
  \end{columns}
\end{frame}

\begin{frame}{Chaîne de traitement}
  \centering
  \begin{tikzpicture}[node distance=0.9cm and 0.8cm,
      box/.style={draw=accent, fill=accent!10, rounded corners, minimum width=2.4cm, minimum height=1cm, align=center, font=\small},
      db/.style={box, cylinder, shape border rotate=90, aspect=0.15, minimum height=1.3cm},
      chk/.style={box, draw=green, fill=green!10}]
    \node[box] (noaa) {Fichiers NOAA\\\footnotesize by\_station, by\_year};
    \node[box, right=of noaa] (dl) {download.py\\\footnotesize manifeste SHA-256};
    \node[box, right=of dl] (load) {load.py\\\footnotesize chargement et F5};
    \node[db, right=of load] (db) {PostgreSQL\\ou Citus};
    \node[box, below=1.2cm of db] (svc) {Services F1--F4\\\footnotesize ligne de commande};
    \node[chk, below=1.2cm of dl] (ver) {verify/\\\footnotesize relit les fichiers};
    \node[chk, below=1.2cm of load] (cmp) {compare.py};
    \draw[-Stealth, thick] (noaa) -- (dl);
    \draw[-Stealth, thick] (dl) -- (load);
    \draw[-Stealth, thick] (load) -- (db);
    \draw[-Stealth, thick] (db) -- (svc);
    \draw[-Stealth, thick, green] (noaa) |- (ver);
    \draw[-Stealth, thick, green] (ver) -- (cmp);
    \draw[-Stealth, thick, green] (svc) -- (cmp);
  \end{tikzpicture}
  \vspace{0.6em}

  \small Un seul chargeur pour le chargement initial et pour F5. La vérification n'utilise ni la base ni le code des services.
\end{frame}

\begin{frame}{Modèle de données}
  \small
  \begin{table}
    \centering
    \begin{tabular}{|l|p{6cm}|l|}
      \hline
      \textbf{Table} & \textbf{Contenu} & \textbf{Sous Citus} \\ \hline
      \code{observations} & une ligne par mesure ; clé (station, élément, date) ; valeur entière en dixièmes, flags, lot & répartie \\ \hline
      \code{monthly\_stats} & par station, élément, mois : comptes, somme, min, max ; tenue à jour par F5 & répartie \\ \hline
      \code{observations\_history} & versions remplacées par F5 & répartie \\ \hline
      \code{stations} & métadonnées, maille de 0,5° & copiée partout \\ \hline
      \code{batches} & lots d'import, fichiers, SHA-256 & copiée partout \\ \hline
    \end{tabular}
  \end{table}
  \begin{block}{Calculs exacts}
    Entiers et \code{numeric}, arrondi à 0,01 seulement à la sortie : les sommes ne dépendent pas de l'ordre des additions, donc local et distribué donnent \textbf{exactement} le même résultat.
  \end{block}
\end{frame}

\begin{frame}{Partie I : deux choix de conception mesurés}
  \small
  \begin{table}
    \centering
    \begin{tabular}{|l|r|r|l|}
      \hline
      \textbf{Mesure (médiane)} & \textbf{Retenu} & \textbf{Alternative} & \textbf{Décision} \\ \hline
      F3 : \code{monthly\_stats} / observations brutes & 0,59 s & 5,81 s & agrégats gardés (\textbf{10$\times$}) \\ \hline
      F1 : une ligne par mesure / par jour & 0,010 s & 0,016 s & \\
      F2 : une ligne par mesure / par jour & 0,10 s & 0,21 s & table large \\
      TMAX $-$ TMIN le même jour & 2,04 s & 0,27 s & écartée \\
      Taille & 1 & 1/3,3 & \\ \hline
    \end{tabular}
  \end{table}
  \textbf{Table large} (une ligne par station et par jour) : 3,3 fois plus petite et 7 fois plus rapide pour comparer TMAX et TMIN, mais F1 et F2, les services du sujet, sont 1,5 à 2 fois plus lents. Résultats identiques vérifiés dans chaque cas.
\end{frame}

%====================================================
\section{Services}

\begin{frame}{Les services et leurs définitions}
  \small
  \begin{table}
    \centering
    \begin{tabular}{|l|p{11cm}|}
      \hline
      \textbf{F1} & Zone = pays et/ou rectangle ; résumé mensuel d'une station : jours attendus, valides, invalides, traces, disponibilité, moyenne ou cumul \\ \hline
      \textbf{F2} & Jours consécutifs avec valeur valide $\geq$ seuil ; un jour absent ou invalide coupe ; décembre--janvier = un seul épisode \\ \hline
      \textbf{F3} & Mois à $\geq$ 80\,\% de jours valides ; normale sur $\geq$ 24 années de 1991--2020 ; moyenne par maille de 0,5° puis entre mailles ($\geq$ 5) \\ \hline
      \textbf{F4} & Inhabituel = au-dessus du quantile 95\,\% de la station pour ce mois ; jour régional = $\geq$ 3 mailles \textbf{et} $\geq$ 25\,\% des mailles qui rapportent \\ \hline
      \textbf{F5} & Écrase et archive ; réimporter ne change rien ; lot contradictoire refusé \\ \hline
    \end{tabular}
  \end{table}
  \textbf{Mailles :} les stations ne sont pas réparties uniformément ; compter des mailles évite qu'une ville dense pèse plus qu'une région entière.
\end{frame}

\begin{frame}[fragile]{F2 : les jours consécutifs en une requête}
  \begin{columns}[T]
    \begin{column}{0.6\textwidth}
\begin{lstlisting}[style=sql]
WITH days AS (  -- jours valides >= seuil
  SELECT station, date FROM observations
  WHERE element = 'TMAX' AND value >= 400
), runs AS (    -- date - rang : constant
  SELECT station, date, date - (row_number()
    OVER (PARTITION BY station ORDER BY date)
  )::int AS grp
  FROM days)
SELECT station, min(date), max(date), count(*)
FROM runs GROUP BY station, grp
HAVING count(*) >= 3
\end{lstlisting}
    \end{column}
    \begin{column}{0.36\textwidth}
      \small
      \begin{tabular}{|l|c|l|}
        \hline
        \textbf{date} & \textbf{rang} & \textbf{date $-$ rang} \\ \hline
        7 juil. & 1 & 6 juil. \\
        8 juil. & 2 & 6 juil. \\
        9 juil. & 3 & 6 juil. \\ \hline
        11 juil. & 4 & \textcolor{red}{7 juil.} \\
        12 juil. & 5 & \textcolor{red}{7 juil.} \\ \hline
      \end{tabular}
      \vspace{0.5em}

      Le 10 manque : nouveau groupe.\\[0.4em]
      Découpé par station, la clé de répartition : chaque worker calcule seul.
    \end{column}
  \end{columns}
\end{frame}

\begin{frame}[fragile]{F5 : un import idempotent}
  \begin{columns}[T]
    \begin{column}{0.56\textwidth}
\begin{lstlisting}[style=sql]
INSERT INTO observations (...)
SELECT DISTINCT ...
FROM observations_staging
ON CONFLICT (station, element, date)
DO UPDATE SET value = excluded.value, ...
WHERE (observations.value, ...)
  IS DISTINCT FROM (excluded.value, ...)
\end{lstlisting}
    \end{column}
    \begin{column}{0.4\textwidth}
      \small
      \begin{enumerate}
        \item \code{COPY} vers une table de chargement
        \item Lot contradictoire ? refusé
        \item Versions remplacées archivées
        \item Écriture du nouveau et du modifié
        \item \code{monthly\_stats} recalculée pour les mois touchés
      \end{enumerate}
      Tout dans une transaction.
    \end{column}
  \end{columns}
  \begin{block}{Vérifié}
    Réimport des 21,46 M lignes : 0 nouvelle, 0 modifiée. Fichier 2025 complet : 697\,993 lignes en 35 s, nouvelles stations ajoutées.
  \end{block}
\end{frame}

\begin{frame}{Traçabilité et vérification}
  \begin{columns}[T]
    \begin{column}{0.48\textwidth}
      \begin{block}{Traçabilité}
        \small
        Chaque résultat contient un champ \code{sources}.\\
        La commande \code{trace} le transforme en liste d'observations : station, date, valeur, flags, lot d'import, donc fichier et SHA-256.
      \end{block}
    \end{column}
    \begin{column}{0.48\textwidth}
      \setbeamercolor{block title}{fg=white,bg=green}
      \begin{block}{Vérification indépendante}
        \small
        Scripts Python qui relisent les fichiers NOAA, sans la base.\\
        \code{verify/compare.py} : 4 cas, \textbf{0 écart}, classements compris (F2 : 220 et 16\,126 épisodes).
      \end{block}
    \end{column}
  \end{columns}
  \vspace{0.5em}
  \small
  \textbf{Exemples :} F3 Allemagne décembre 2015 +5,16 °C ; F2 SPE00120503, 11 jours $\geq$ 40 °C en juillet 2022 ; F4 Allemagne 166 événements, le plus sévère du 1\textsuperscript{er} au 12 novembre 2015.
\end{frame}

%====================================================
\section{Distribution}

\begin{frame}{Répartition par hachage de la station}
  \begin{columns}[c]
    \begin{column}{0.48\textwidth}
      \centering
      \begin{tikzpicture}[node distance=0.8cm and 0.1cm,
          box/.style={draw=accent, fill=accent!10, rounded corners, minimum width=2cm, minimum height=0.9cm, align=center, font=\small}]
        \node[box] (app) {Application};
        \node[box, below=of app] (co) {Coordinateur};
        \node[box, below left=of co] (w1) {Worker 1\\\footnotesize 16 shards};
        \node[box, below right=of co] (w2) {Worker 2\\\footnotesize 16 shards};
        \draw[-Stealth, thick] (app) -- (co);
        \draw[-Stealth, thick] (co) -- (w1);
        \draw[-Stealth, thick] (co) -- (w2);
      \end{tikzpicture}
    \end{column}
    \begin{column}{0.5\textwidth}
      \small
      \begin{itemize}
        \item Une station = un worker : \textbf{F1 contacte un seul worker} (\code{Task Count: 1})
        \item F2 et F5 calculés entièrement sur chaque worker
        \item F3, F4 : 32 tâches, fusion par le coordinateur
        \item Tables liées colocalisées ; \code{stations}, \code{batches} copiées partout
        \item Écartés : par date (épisodes coupés), par pays (2 valeurs)
      \end{itemize}
    \end{column}
  \end{columns}
\end{frame}

\begin{frame}{Comparaison et panne}
  \begin{block}{Cluster de test (1 coordinateur + 2 workers, 65 stations, 1,8 M observations)}
    \begin{itemize}
      \item 12 commandes : sortie \textbf{identique octet pour octet} au local
      \item Worker arrêté : F1 sur une station de l'autre worker répond ; F3 affiche « ÉCHEC », aucun résultat partiel, code 2
    \end{itemize}
  \end{block}
  \setbeamercolor{block title}{fg=white,bg=red!80!black}
  \begin{block}{Données complètes}
    \afaire{déploiement 3 workers, comparaison automatique sur les 21,46 M lignes, répartition par worker.}
  \end{block}
\end{frame}

%====================================================
\section{Expériences}

\begin{frame}{E1 et E2}
  \begin{block}{E1 -- Volume}
    \afaire{3 tailles (7,14 M, intermédiaire, 21,46 M) : temps de chargement, taille, F1 et F3.}
  \end{block}
  \begin{block}{E2 -- Local contre distribué}
    \afaire{local, Citus-2, Citus-3 sur F1 (une station) et F2/F3 (un pays).}
  \end{block}
\end{frame}

\begin{frame}{E3 -- Charge variable : prédiction déposée avant les mesures}
  \small
  Charge de 200 requêtes : routine (F1, une station) et étude régionale (F2/F3, un pays), de 0 à 100\,\% de régional.
  \begin{enumerate}
    \item \textbf{Routine seule :} le local gagne (surcoût fixe du coordinateur, 1,3 à 2$\times$)
    \item \textbf{Régional seul :} Citus-2 gagne grâce à F2 (32 tâches en parallèle, 1,5 à 3$\times$)
    \item \textbf{Mélange :} bascule tôt, vers 1 à 5\,\% de régional
  \end{enumerate}
  \textbf{Risques :} mémoire partagée, un seul poste, plan d'exécution local.
  \begin{block}{Mesures}
    \afaire{résultats et confrontation avec la prédiction.}
  \end{block}
\end{frame}

\begin{frame}{E4 -- Limite}
  \begin{block}{Piste identifiée}
    F4 recalcule les quantiles de référence (30 ans) à chaque requête : 29 s sur l'Allemagne.\\
    Une table de seuils par station et mois, tenue à jour par F5, les rendrait instantanés.
  \end{block}
  \afaire{mesure si réalisée.}
\end{frame}

%====================================================
\section{Démonstration et limites}

\begin{frame}[fragile]{Démonstration (5 minutes)}
  \small
  \begin{enumerate}
    \item \textbf{F1} : chercher, puis résumer une station
\begin{lstlisting}[basicstyle=\ttfamily\scriptsize,frame=none]
ws stations -z SP --name madrid
ws station SP000003195 --start 2022-07-01 --end 2022-07-31
\end{lstlisting}
    \item \textbf{F2, F3} : épisodes de chaleur, mois anormaux
\begin{lstlisting}[basicstyle=\ttfamily\scriptsize,frame=none]
ws episodes -z SP --start 2015-01-01 --end 2024-12-31 --threshold 40
ws anomalies -z GM -z SP --element TMAX
\end{lstlisting}
    \item \textbf{Traçabilité} : les jours d'un épisode
\begin{lstlisting}[basicstyle=\ttfamily\scriptsize,frame=none]
ws trace --station SPE00120503 --element TMAX --start 2022-07-07 --end 2022-07-17
\end{lstlisting}
    \item \textbf{F5} : importer un lot, relancer F3 ; \textbf{Citus} : \code{-{}-explain}
  \end{enumerate}
  \code{ws} = \code{uv run python -m weatherscope.cli}
\end{frame}

\begin{frame}{Limites}
  \begin{itemize}
    \item Un seul poste : plusieurs conteneurs ne valent pas un vrai cluster de machines
    \item Pas de réplication : une panne de worker bloque les requêtes qui en ont besoin
    \item F4 lent (quantiles recalculés)
    \item Métadonnées NOAA laissées telles quelles (longitudes fausses)
    \item Interface en ligne de commande seulement
  \end{itemize}
  \vspace{0.5em}
  \small Toutes les décisions, leurs raisons et les alternatives : \code{docs/decisions.md} (D1--D26).
\end{frame}

%====================================================
\begin{frame}{Merci}
  \centering\Large Merci pour votre attention !\\[1em]
  \normalsize Questions et défi en direct
\end{frame}

\end{document}
